% !TEX root = ../../ctfp-print.tex

\lettrine[lhang=0.17]{在}{本书的第一部分}我论证了范畴论和编程都关注可组合性。在编程中，你持续分解问题直到达到可处理的细节层级，依次解决每个子问题，然后自底向上重组解决方案。粗略来说有两种实现方式：告诉计算机做什么，或者告诉它怎么做。前者称为声明式，后者称为命令式。

即使在最基本层面也能看到这种区别。组合本身可以声明式地定义；例如，\code{h}是\code{g}在\code{f}之后的复合：

\src{snippet01}
或者命令式地定义；例如，先调用\code{f}，记住该调用结果，然后用该结果调用\code{g}：

\src{snippet02}
程序的命令式版本通常被描述为按时间顺序排列的一系列动作。特别地，\code{g}的调用不能早于\code{f}的执行完成。至少这是概念上的画面——在惰性语言中，通过\emph{按需调用}的传参机制，实际执行可能有所不同。

事实上，取决于编译器的智能程度，声明式和命令式代码的执行方式可能几乎没有差异。但两种方法论在解决问题的方式上，以及生成代码的可维护性和可测试性方面，有时存在显著差异。

核心问题是：当面对一个问题时，我们是否总能在声明式和命令式方法之间进行选择？如果存在声明式解法，它是否总能转化为计算机代码？这个问题的答案远非显而易见，如果我们能找到答案，可能会彻底革新我们对宇宙的理解。

\begin{wrapfigure}{R}{0pt}
  \includegraphics[width=0.5\textwidth]{images/asteroids.png}
\end{wrapfigure}

让我展开说明。物理学中也存在类似的二元性，这要么指向某种深层的基本原理，要么揭示了我们思维运作的方式。理查德·费曼在其量子电动力学研究中就将这种二元性视为灵感来源。

大多数物理定律有两种表述形式。一种采用局部（或无穷小）考虑。我们观察系统在某个小邻域的状态，预测其在下一时刻的演化。这通常通过微分方程表达，需要在一段时间内积分或累加。

注意这种方法如何类似于命令式思维：我们通过遵循一系列依赖前一步结果的小步骤到达最终解。事实上，物理系统的计算机模拟通常是将微分方程转化为差分方程并迭代求解。这就是《太空陨石》游戏中飞船动画的实现方式。在每个时间步长，飞船位置通过添加一个小增量而改变，该增量等于其速度乘以时间差。速度本身则通过与加速度成比例的小增量改变，而加速度由力除以质量给出。

这些是对应牛顿运动定律的微分方程的直接编码：
\begin{align*}
  F & = m \frac{dv}{dt} \\
  v & = \frac{dx}{dt}
\end{align*}
类似方法可应用于更复杂的问题，例如使用麦克斯韦方程组模拟电磁场传播，甚至使用格点量子色动力学（\acronym{QCD}）模拟质子内部夸克和胶子的行为。

这种受数字计算机鼓励的时空离散化与局部思维结合，在史蒂芬·沃尔夫勒姆试图将整个宇宙的复杂性归约为细胞自动机系统的壮举中达到了极致。

另一种方法是全局性的。我们观察系统的初始和终态，通过最小化某个泛函来计算连接它们的轨迹。最简单的例子是费马的最短时间原理。它指出光沿使传播时间最小化的路径行进。特别地，在没有反射或折射物体的情况下，从点$A$到点$B$的光线会选择最短路径——即直线。但光在致密（透明）材料（如水或玻璃）中传播较慢。因此如果你选择起点在空气中，终点在水下，光线更倾向于在空气中走更长距离，然后通过水路"抄近道"。这种最短时间路径使光线在空气-水界面发生折射，产生斯涅尔折射定律：
\begin{equation*}
  \frac{\sin(\theta_1)}{\sin(\theta_2)} = \frac{v_1}{v_2}
\end{equation*}
其中$v_1$是光在空气中的速度，$v_2$是光在水中的速度。

\begin{figure}[H]
  \centering
  \includegraphics[width=0.3\textwidth]{images/snell.jpg}
\end{figure}

\noindent
全部经典力学都可以从最小作用量原理导出。作用量可通过积分拉格朗日量计算，而拉格朗日量是动能与势能之差（注意：是差值而非和值——和值才是总能量）。当你发射迫击炮弹命中目标时，弹丸会先上升到重力势能更高的高度，在那里积累负的作用量贡献。它也会在抛物线顶端减速以最小化动能，然后加速穿过低势能区域。

\begin{figure}[H]
  \centering
  \includegraphics[width=0.35\textwidth]{images/mortar.jpg}
\end{figure}

\noindent
费曼的最大贡献在于认识到最小作用量原理可以推广到量子力学。在那里同样将问题表述为初始态和终态的关系。通过计算两态间的费曼路径积分来确定跃迁概率。

\begin{figure}[H]
  \centering
  \includegraphics[width=0.35\textwidth]{images/feynman.jpg}
\end{figure}

\noindent
关键在于我们描述物理定律的方式存在一种神秘的未解二元性。我们可以使用事物按小增量顺序发生的局部图景，也可以使用声明初始和终态后中间一切自然呈现的全局图景。

全局方法同样可用于编程，例如实现光线追踪时。我们声明眼睛位置和光源位置，计算光线可能采取的连接路径。虽然没有显式最小化每条光线的飞行时间，但我们确实通过斯涅尔定律和反射几何达到相同效果。

局部与全局方法最大的差异在于对空间特别是时间的处理。局部方法拥抱即时满足的当下性，而全局方法采取长期静态视角，仿佛未来已被预先注定，我们只是在分析某个永恒宇宙的属性。

这一点在函数式响应编程（\acronym{FRP}）处理用户交互时体现得尤为明显。与其为每个可能的用户操作编写独立的处理器并共享可变状态，\acronym{FRP}将外部事件视为无限列表并对之应用变换序列。从概念上说，所有未来操作的列表作为输入数据已存在于程序中。从程序视角看，$\pi$的小数位列表、伪随机数列表或鼠标位置列表并无区别。在任一情况下，若要获取第$n$项，必须先遍历前$n-1$项。当应用于时序事件时，我们称此性质为\emph{因果性}。

这与范畴论有何关联？我认为范畴论鼓励全局视角因而支持声明式编程。首先，不同于微积分，它没有内置的距离、邻域或时间概念。我们拥有的只是抽象对象及其抽象连接。若你能通过一系列步骤从$A$到达$B$，你就能单步到达。更重要的是，范畴论的主要工具是万有构造——全局方法的典范。我们在范畴积定义中见过其实例：通过指定其属性完成定义——典型的声明式方法。它是配备两个投影的对象，且是此类对象中最优者——它优化了其他同类对象投影分解的性质。

\begin{figure}[H]
  \centering
  \includegraphics[width=0.35\textwidth]{images/productranking.jpg}
\end{figure}

\noindent
将此与费马最短时间原理或最小作用量原理对比。

反之，对比传统笛卡尔积定义——更具命令性。你通过从一个集合选元素、另一个集合选元素来创建积元素。这是一个创建有序对的操作指南。还有一个拆解有序对的操作指南。

几乎所有编程语言（包括Haskell这类函数式语言）都将乘积类型、余积类型和函数类型作为内置结构，而非通过万有构造定义；尽管已有尝试创建范畴论编程语言（参见，例如，
\urlref{http://web.sfc.keio.ac.jp/~hagino/thesis.pdf}{萩野达也的博士论文}）。

无论是否直接使用，范畴论定义都在证明既有编程结构合理性的同时催生新结构。最重要的是，范畴论提供了在声明层面对程序推理的元语言。它还鼓励在编写代码前先进行问题规范的推理。